\uuid{7Wzz}
\titre{Noyau/Image/Injection/Surjection/Bijection/Théorème du rang}
\auteur{hanani ; mzari}
\organisation{exo7}
\licence{CC BY-NC-SA 4.0}
\datecreate{2020-01-20}
\niveau{L1}
\module{Algèbre}
\chapitre{Application linéaire}
\sousChapitre{Définition}
\difficulte{2}
\theme{Noyau et image, Noyau, Image, Injection, Surjection, Bijection, Théorème du rang}
\exo7qcmid{lille-2.applications-lineaires.017}

\contenu{
\begin{question}
\qclassification{107.01}
\qtags{motcle=Noyau/Image/Injection/Surjection/Bijection/Théorème du rang}
On considère l' application linéaire :
$$\begin{array}{rccc}f:&\Rr^3&\to&\Rr^2\\
& (x,y,z)&\to &(x+y+z, x+y-z). \,  \end{array}$$
Quelles sont les assertions vraies ?
\begin{answers}
\good{$\dim \ker f = 1$}.
\bad{$f$ est injective}.
\bad{$\mbox{rg} (f) =1$}.
\good{$f$ n'est pas bijective}.
\end{answers}
\begin{explanations}  $\ker f = \{(x,y,z) \in \Rr^3 \; ; \; (x+y+z,x+y-z)=(0,0,0)\} =\{(x,-x,0)  \; ; \; x\in \Rr\}$. Donc $ \{(1,-1,0)\}$ est une base de $\ker f$. Comme $\ker f \neq \{(0,0,0)\}$, $f$ n'est pas injective, donc $f$ n'est pas bijective. Du théorème du rang, on déduit que $\mbox{rg} (f)= \dim \Im f = 2$.
\end{explanations}
\end{question}
}
